\section{The softmax gate, its likelihood, EM and backpropagation}

This section treats the textbook gate (Bishop's softmax) together with MLP experts. Section 5 then turns the gate into an HMM. The machinery built here (complete-data likelihood, EM, the gradient identity, backpropagation) carries over unchanged.

\subsection{Where the softmax gate comes from}

\paragraph{Derivation 1: Bayes' rule with Gaussian class-conditional densities (Bishop \S4.2.2).} Suppose the gate input $u$ is generated by the regime: first draw $k$ with probability $p(k)$, then $u\sim\N(u;m_k,\Sigma)$ with a covariance $\Sigma$ shared by all regimes. Bayes' rule gives
\begin{equation}
p(k\mid u)=\frac{p(u\mid k)\,p(k)}{\sum_j p(u\mid j)\,p(j)}=\frac{\exp(a_k)}{\sum_j\exp(a_j)},\qquad a_k=\log\big[p(u\mid k)\,p(k)\big].
\end{equation}
Any posterior can be written as a softmax of log-joint probabilities, so this step is only a rewriting. The assumption enters when we expand the Gaussian:
\begin{equation}
a_k=-\tfrac12 u\T\Sigma^{-1}u+m_k\T\Sigma^{-1}u-\tfrac12 m_k\T\Sigma^{-1}m_k+\log p(k)+\text{const}.
\end{equation}
The quadratic term $-\frac12 u\T\Sigma^{-1}u$ and the constant are the same for every $k$, so they cancel between numerator and denominator. What remains is linear in $u$:
\begin{equation}
p(k\mid u)=\softmax_k\big(w_k\T u+w_{k0}\big),\qquad w_k=\Sigma^{-1}m_k,\qquad w_{k0}=-\tfrac12m_k\T\Sigma^{-1}m_k+\log p(k).
\end{equation}
\textbf{A softmax of linear functions is therefore the exact posterior of a generative model of the gate input.} If the covariances differ across regimes, the quadratic terms no longer cancel and the logits become quadratic in $u$. More generally, any exponential-family class-conditional density with a shared dispersion parameter gives linear logits (Bishop \S4.2.4).

\paragraph{Derivation 2: linear log-odds (the multinomial logit).} Assume directly that the log-odds of regime $k$ against a reference regime $K$ are linear in the input:
\begin{equation}
\log\frac{p(k\mid u)}{p(K\mid u)}=w_k\T u\qquad(k=1,\dots,K-1).
\end{equation}
Exponentiate, impose $\sum_k p(k\mid u)=1$, and set $w_K=0$:
\begin{equation}
p(k\mid u)=\frac{\exp(w_k\T u)}{1+\sum_{j<K}\exp(w_j\T u)}=\frac{\exp(w_k\T u)}{\sum_{j=1}^K\exp(w_j\T u)}.
\end{equation}
This is the canonical-link generalised linear model for a categorical outcome.

\paragraph{Identifiability.} Adding the same vector $c$ to every $w_k$ multiplies the numerator and the denominator by $\exp(c\T u)$, so $\pi$ is unchanged. Only $K-1$ of the $w_k$ can be identified, which is why one of them is fixed at zero.

\paragraph{Why it is called ``soft max''.} With a temperature $\tau$, $\softmax_k(a/\tau)$ tends to the one-hot vector of $\argmax_k a_k$ as $\tau\to0$ and to the uniform distribution as $\tau\to\infty$. Its derivative, which we use repeatedly below, is
\begin{equation}\label{eq:softmax-deriv}
\frac{\partial\pi_j}{\partial a_k}=\pi_j\,(\delta_{jk}-\pi_k).
\end{equation}

\begin{keybox}{The bridge to Section 5}
The HMM filter of Section 5 is also a softmax. Its logits are the log of the emission density plus the log of the prior that the transition matrix carries forward from yesterday:
\[
\xi_{t\mid t}(k)=\softmax_k\Big(\log\eta_t(k)+\log\textstyle\sum_jA_{jk}\,\xi_{t-1\mid t-1}(j)\Big).
\]
In the softmax gate, the ``prior'' part of the logit is a fixed bias $w_{k0}$. In the HMM it is recomputed every day from yesterday's belief. That is the precise sense in which the HMM gate is a softmax gate with memory.
\end{keybox}

\subsection{The generative model of the softmax-gated MLP MoE}

The data are $\{(u_n,x_n,y_n)\}_{n=1}^N$: $u_n$ is the gate input (in Bishop $u_n=x_n$), $x_n$ the expert input and $y_n$ the target.
\begin{align}
z_n\mid u_n&\sim\text{Cat}\big(\pi_1(u_n;w),\dots,\pi_K(u_n;w)\big),\qquad \pi_k(u;w)=\softmax_k\big(w_k\T u\big),\\
y_n\mid x_n,z_{nk}=1&\sim\N\big(\mu_k(x_n;\theta_k),\,s_k^2\big),\qquad \mu_k\text{ an MLP (Section 3)}.
\end{align}

\begin{center}
\begin{tikzpicture}[node distance=14mm]
\node[obs] (u) {$u_n$};
\node[lat,right=of u] (z) {$z_n$};
\node[obs,right=of z] (y) {$y_n$};
\node[obs,above=9mm of y] (x) {$x_n$};
\draw[edge] (u)--(z); \draw[edge] (z)--(y); \draw[edge] (x)--(y);
\node[par,below=8mm of z,label=below:{$w$}] (w) {}; \draw[edge] (w)--(z);
\node[par,below=8mm of y,label=below:{$\theta_k,s_k$}] (th) {}; \draw[edge] (th)--(y);
\begin{scope}[on background layer]
\node[plate,fit=(u)(z)(y)(x),label={[anchor=north west]south west:$n=1,\dots,N$}] {};
\end{scope}
\end{tikzpicture}
\end{center}

\paragraph{Complete-data likelihood.} With one-hot $z_n$,
\begin{equation}
p(Y,Z\mid X,U;\Psi)=\prod_{n=1}^N\prod_{k=1}^K\Big[\pi_k(u_n;w)\,\N\big(y_n;\mu_k(x_n;\theta_k),s_k^2\big)\Big]^{z_{nk}}.
\end{equation}
\paragraph{Observed-data log-likelihood.} Summing out each $z_n$ (Appendix A(c) shows the step) gives
\begin{equation}\label{eq:softmax-ll}
\ell(\Psi)=\log p(Y\mid X,U;\Psi)=\sum_{n=1}^N\log\sum_{k=1}^K\pi_k(u_n;w)\,\N\big(y_n;\mu_k(x_n;\theta_k),s_k^2\big).
\end{equation}
The sum inside the logarithm couples all the parameters, so there is no closed-form maximiser even when the experts are linear.

\subsection{Maximising it with EM}

\paragraph{E-step.} Using the current parameters $\Psi^{(q)}$, compute the responsibilities
\begin{equation}
\gamma_{nk}=p(z_{nk}=1\mid y_n,x_n,u_n;\Psi^{(q)})=\frac{\pi_k(u_n)\,\N\big(y_n;\mu_k(x_n),s_k^2\big)}{\sum_j\pi_j(u_n)\,\N\big(y_n;\mu_j(x_n),s_j^2\big)}.
\end{equation}
\paragraph{The $Q$-function} is the complete-data log-likelihood with each $z_{nk}$ replaced by $\gamma_{nk}$. It separates into a gate part and one part per expert:
\begin{equation}
Q(\Psi;\Psi^{(q)})=\underbrace{\sum_n\sum_k\gamma_{nk}\log\pi_k(u_n;w)}_{Q_{\text{gate}}(w)}+\sum_k\underbrace{\sum_n\gamma_{nk}\Big[-\tfrac12\log(2\pi s_k^2)-\frac{(y_n-\mu_k(x_n;\theta_k))^2}{2s_k^2}\Big]}_{Q_k(\theta_k,s_k)}.
\end{equation}
\paragraph{M-step.}
\begin{itemize}
  \item \emph{Gate.} $Q_{\text{gate}}$ is a cross-entropy with soft targets $\gamma_{nk}$, i.e.\ a weighted multinomial logistic regression. It is concave in $w$ and is solved by IRLS (Appendix A(g)).
  \item \emph{Experts.} For fixed $s_k$, maximising $Q_k$ over $\theta_k$ is \emph{weighted least squares}, $\min_{\theta_k}\sum_n\gamma_{nk}(y_n-\mu_k(x_n;\theta_k))^2$. For linear experts this has a closed form. For MLPs it does not, so one takes a few gradient steps. Doing so still increases $Q$, which is all EM needs: this is \emph{generalised EM}.
  \item \emph{Noise.} Setting $\partial Q_k/\partial s_k^2=0$ gives $s_k^2=\sum_n\gamma_{nk}(y_n-\mu_k(x_n))^2\big/\sum_n\gamma_{nk}$.
\end{itemize}

\subsection{Why EM never decreases the likelihood}

For \emph{any} distribution $q(Z)$ over the latent variables,
\begin{align}
\log p(Y\mid\Psi)&=\sum_Z q(Z)\log p(Y\mid\Psi)=\sum_Zq(Z)\log\frac{p(Y,Z\mid\Psi)}{p(Z\mid Y,\Psi)}\nonumber\\
&=\underbrace{\sum_Zq(Z)\log\frac{p(Y,Z\mid\Psi)}{q(Z)}}_{\mathcal L(q,\Psi)}\;+\;\underbrace{\sum_Zq(Z)\log\frac{q(Z)}{p(Z\mid Y,\Psi)}}_{\mathrm{KL}(q\,\|\,p(Z\mid Y,\Psi))\;\ge\;0}.
\end{align}
So $\mathcal L(q,\Psi)$ is a lower bound on the log-likelihood (Bishop \S9.4).
\begin{itemize}
  \item The \textbf{E-step} sets $q(Z)=p(Z\mid Y,\Psi^{(q)})$. The KL term becomes zero and the bound touches the log-likelihood at $\Psi^{(q)}$.
  \item With this $q$, $\mathcal L(q,\Psi)=Q(\Psi;\Psi^{(q)})-\sum_Zq\log q$. The second term does not depend on $\Psi$, so the \textbf{M-step}, which maximises $Q$, maximises the bound.
\end{itemize}
Therefore
\begin{equation}
\log p(Y\mid\Psi^{(q+1)})\;\ge\;\mathcal L(q,\Psi^{(q+1)})\;\ge\;\mathcal L(q,\Psi^{(q)})\;=\;\log p(Y\mid\Psi^{(q)}).
\end{equation}
Any $\Psi^{(q+1)}$ that merely \emph{increases} $Q$ keeps this chain of inequalities, which is why generalised EM still works.

\subsection{Why gradient training of the likelihood behaves like EM}

\paragraph{Fisher's identity.} The gradient of the observed-data log-likelihood equals the posterior expectation of the gradient of the complete-data log-likelihood:
\begin{equation}
\nabla_\Psi\log p(Y\mid\Psi)=\frac{\sum_Z\nabla_\Psi p(Y,Z\mid\Psi)}{p(Y\mid\Psi)}=\sum_Z\frac{p(Y,Z\mid\Psi)}{p(Y\mid\Psi)}\nabla_\Psi\log p(Y,Z\mid\Psi)=\E_{p(Z\mid Y,\Psi)}\big[\nabla_\Psi\log p(Y,Z\mid\Psi)\big].
\end{equation}
The right-hand side is exactly $\nabla_\Psi Q(\Psi;\Psi')$ evaluated at $\Psi'=\Psi$. \textbf{A gradient step on the log-likelihood is therefore an E-step (compute $\gamma$ with the current parameters) followed by one small step on the M-step objective.} EM and gradient ascent climb the same function and stop at the same kind of point, where the gradient is zero. EM jumps to the maximum of $Q$ when it can; gradient ascent takes small steps. With MLP experts we cannot jump anyway.

\paragraph{The gradients for one observation.} Write $\ell_n=\log\sum_k\pi_k\N_k$ with $\N_k=\N(y_n;\mu_k,s_k^2)$.
\begin{itemize}
  \item \emph{Expert mean.} Only the $k$-th term depends on $\mu_k$, and $\partial\log\N_k/\partial\mu_k=(y_n-\mu_k)/s_k^2$, so
  \begin{equation}\label{eq:grad-mu}
  \frac{\partial\ell_n}{\partial\mu_k}=\frac{\pi_k\N_k}{\sum_j\pi_j\N_j}\cdot\frac{y_n-\mu_k}{s_k^2}=\gamma_{nk}\,\frac{y_n-\mu_k}{s_k^2}.
  \end{equation}
  \item \emph{Expert noise.} Since $\partial\log\N_k/\partial\log s_k=-1+(y_n-\mu_k)^2/s_k^2$,
  \begin{equation}
  \frac{\partial\ell_n}{\partial\log s_k}=\gamma_{nk}\Big[\frac{(y_n-\mu_k)^2}{s_k^2}-1\Big].
  \end{equation}
  \item \emph{Gate logits} $a_k$, with $\pi=\softmax(a)$. Using \eqref{eq:softmax-deriv},
  \begin{equation}
  \frac{\partial\ell_n}{\partial a_k}=\sum_j\frac{\N_j}{\sum_l\pi_l\N_l}\,\pi_j(\delta_{jk}-\pi_k)=\gamma_{nk}-\pi_k\sum_j\gamma_{nj}=\gamma_{nk}-\pi_k.
  \end{equation}
\end{itemize}
Each gradient is the EM weighting in disguise. The expert that ``owns'' an observation (large $\gamma_{nk}$) receives its error signal, scaled by $1/s_k^2$. The gate is pulled towards the posterior $\gamma$, as in Weigend et al.\ (1995) eqs.~24--25. In our base case the gate is frozen, so only \eqref{eq:grad-mu} and the noise gradient are used.

\subsection{Backpropagation in depth}

Backpropagation is the chain rule applied efficiently to the expert MLP of Section 3, eqs.~\eqref{eq:pre}--\eqref{eq:out}. We minimise the loss $L_n=-\ell_n$ for one observation. Its derivative with respect to expert $k$'s output is, from \eqref{eq:grad-mu},
\begin{equation}
\delta^{\text{out}}_k:=\frac{\partial L_n}{\partial\mu_k}=-\gamma_{nk}\,\frac{y_n-\mu_k}{s_k^2}.
\end{equation}
Drop the index $k$ below: the same steps run inside every expert.

\paragraph{Step 1: forward pass.} Compute and \emph{store} $a^{(\ell)}$ and $h^{(\ell)}$ for $\ell=1,\dots,L$, then the output $\mu$, the responsibility $\gamma$ and $\delta^{\text{out}}$.

\paragraph{Step 2: output layer.} Since $\mu=w^{(L+1)\top}h^{(L)}+b^{(L+1)}$,
\begin{equation}
\frac{\partial L_n}{\partial w^{(L+1)}}=\delta^{\text{out}}\,h^{(L)},\qquad \frac{\partial L_n}{\partial b^{(L+1)}}=\delta^{\text{out}}.
\end{equation}

\paragraph{Step 3: define the error at every hidden layer,} $\delta^{(\ell)}:=\partial L_n/\partial a^{(\ell)}\in\R^{H_\ell}$. For the last hidden layer, $\mu$ depends on $a^{(L)}_m$ only through $h^{(L)}_m=g(a^{(L)}_m)$, so
\begin{equation}
\delta^{(L)}_m=\delta^{\text{out}}\,w^{(L+1)}_m\,g'\big(a^{(L)}_m\big),\qquad\text{i.e.}\qquad\delta^{(L)}=g'\big(a^{(L)}\big)\odot w^{(L+1)}\,\delta^{\text{out}}.
\end{equation}

\paragraph{Step 4: the backward recursion.} Layer $\ell$ affects the loss only through the pre-activations of layer $\ell+1$, $a^{(\ell+1)}_j=\sum_mW^{(\ell+1)}_{jm}\,g(a^{(\ell)}_m)+b^{(\ell+1)}_j$, so $\partial a^{(\ell+1)}_j/\partial a^{(\ell)}_m=W^{(\ell+1)}_{jm}\,g'(a^{(\ell)}_m)$. The chain rule sums over every path from $a^{(\ell)}_m$ to the loss:
\begin{equation}
\delta^{(\ell)}_m=\sum_{j=1}^{H_{\ell+1}}\delta^{(\ell+1)}_j\,W^{(\ell+1)}_{jm}\,g'\big(a^{(\ell)}_m\big)\qquad\Longleftrightarrow\qquad\delta^{(\ell)}=g'\big(a^{(\ell)}\big)\odot\Big(W^{(\ell+1)\top}\delta^{(\ell+1)}\Big).
\end{equation}

\paragraph{Step 5: weight gradients.} Since $a^{(\ell)}_j=\sum_mW^{(\ell)}_{jm}h^{(\ell-1)}_m+b^{(\ell)}_j$,
\begin{equation}
\frac{\partial L_n}{\partial W^{(\ell)}_{jm}}=\delta^{(\ell)}_j\,h^{(\ell-1)}_m\quad\Longleftrightarrow\quad\frac{\partial L_n}{\partial W^{(\ell)}}=\delta^{(\ell)}\,h^{(\ell-1)\top},\qquad\frac{\partial L_n}{\partial b^{(\ell)}}=\delta^{(\ell)}.
\end{equation}
Each gradient is ``error arriving at the unit'' times ``signal that entered the unit''.

\paragraph{A tiny example.} One hidden unit, squared-error loss: $f(x)=w_2\,g(w_1x+b_1)+b_2$, $a=w_1x+b_1$, $L=\frac12(y-f)^2$.
\begin{equation}
\frac{\partial L}{\partial w_2}=-(y-f)\,g(a),\quad\frac{\partial L}{\partial b_2}=-(y-f),\quad\frac{\partial L}{\partial w_1}=-(y-f)\,w_2\,g'(a)\,x,\quad\frac{\partial L}{\partial b_1}=-(y-f)\,w_2\,g'(a).
\end{equation}
The error $-(y-f)$ is computed once and reused, travelling backwards through $w_2$ and $g'$. That reuse is the whole trick.

\paragraph{Why this is efficient.} One backward pass costs about as much as one forward pass and returns the gradient with respect to \emph{all} parameters. Finite differences would need one extra forward pass per parameter. Reverse-mode automatic differentiation (PyTorch's autograd) carries out exactly Steps 1--5 for any graph we write down, including the mixture log-likelihood and the gate.

\paragraph{The update.} Average the gradients over a mini-batch $\mathcal B$ and step downhill:
\begin{equation}
\theta\leftarrow\theta-\eta\,\frac{1}{|\mathcal B|}\sum_{n\in\mathcal B}\nabla_\theta L_n.
\end{equation}
Adam, which the code uses, rescales each coordinate of this step by running estimates of the gradient's mean and root mean square. It is still a descent step on the same negative log-likelihood.

\paragraph{Three consequences for our base case.}
\begin{enumerate}
  \item \emph{Frozen base.} $f_0$ appears inside every $\mu_k$, but its parameters $\phi$ are excluded from the update, so the error stops at $f_0$.
  \item \emph{Zero-initialised final layer.} At step zero $w^{(L+1)}=0$, so $\delta^{(L)}=0$ and the hidden layers receive no gradient. They start learning one step later, once the final layer has moved.
  \item \emph{Mini-batches under the per-date likelihood.} The per-date likelihood (Section 7) couples all stocks of a date through the shared regime, so $\gamma$ must be computed over \emph{whole dates}. Mini-batches must therefore consist of complete dates. Under the per-row likelihood, which the study uses (Q24), any set of rows will do.
\end{enumerate}
\newpage
